% Source: https://artofproblemsolving.com/community/c3h1430822_trapezoid_question__find_the_angle
% Stable identity: aops-topic-1430822-q01
Take $RS=SP=1$, so $PQ=3$. Mark $T$ on $PQ$ with $PT=1$.

\begin{asy}
usepackage("amsmath");
size(6cm);
defaultpen(fontsize(10pt));
pair ft(pair P,pair U,pair V){return U+dot(P-U,V-U)/dot(V-U,V-U)*(V-U);}
path ra(pair P,pair u,pair v,real s=0.2){return P+s*unit(u)--P+s*unit(u)+s*unit(v)--P+s*unit(v);}
pair P=(0,0),Q=(3,0),Sp=(0.5,sqrt(3)/2),R=(1.5,sqrt(3)/2),T=(1,0);
draw(P--Q--R--Sp--cycle,linewidth(0.9));
draw(T--R,red+dashed);
draw(ra(R,T-R,Q-R,0.13),red);
draw(arc(Sp,0.18,-120,0)); label("$120^\circ$",Sp+0.2*dir(-60),SE,fontsize(8pt));
draw(arc(Q,0.45,150,180),blue); label("$30^\circ$",Q+0.6*dir(167),W,blue+fontsize(8pt));
dot(P);dot(Q);dot(R);dot(Sp);dot(T);
label("$P$",P,SW); label("$Q$",Q,SE); label("$R$",R,N); label("$S$",Sp,N); label("$T$",T,S);
label("$1$",(Sp+R)/2,N); label("$1$",(P+Sp)/2,NW); label("$1$",(P+T)/2,S); label("$2$",(T+Q)/2,S);
label("\mbox{}",(min(currentpicture,user=true)+max(currentpicture,user=true))/2,black); // keeps the text after the figure black when LaTeX typesets the labels inline
\end{asy}

Since $PT\parallel SR$ and $PT=SR$, $PTRS$ is a parallelogram (in fact a rhombus), so $TR=SP=1$ and $TR\parallel PS$. Co-interior angles give $\angle SPQ=180^\circ-120^\circ=60^\circ$, and then corresponding angles give $\angle RTQ=\angle SPQ=60^\circ$.

In triangle $TRQ$ we have $TR=1$, $TQ=3-1=2$ and $\angle RTQ=60^\circ$. By the law of cosines,
\[
RQ^2=1^2+2^2-2\cdot1\cdot2\cos60^\circ=3 .
\]
Then $TR^2+RQ^2=1+3=4=TQ^2$, so $\angle TRQ=90^\circ$, and triangle $TRQ$ is a $30$-$60$-$90$ triangle with hypotenuse $TQ=2$ and shorter leg $TR=1$. The angle opposite the shorter leg is
\[
\angle PQR=\angle TQR=\boxed{\textbf{(D)}\ 30^\circ}.
\]
